\chapter{Introduction}

\section{Overview}

Zinc is a statically-typed, compiled systems programming language designed for clarity, performance, and expressiveness.
It compiles native binaries via LLVM and supports C interoperability through \zn{foreign} declarations.

\section{Compiler Pipeline}

The Zinc compiler is a multi-phase pipeline written in C. The phases run sequentially on each source file:

\begin{center}
\begin{tabular}{lll}
\toprule
\textbf{Phase} & \textbf{Source file} & \textbf{Responsibility}\\
\midrule

Lexical analysis & \texttt{zlenx.c} & Tokenise source test into a stream od tokens.\\

Syntax analysis & \texttt{zparse.c} & Build the AST from the token stream.\\

Semantic analysis & \texttt{zsem.c} & Resolve types and symbol resolution\\

Code generation & \texttt{zgen.c} & Emit LLVM IR and link with lld.
\bottomrule
\end{tabular}
\end{center}

Supporting components:

\begin{description}
    \item[\texttt{zmacro.c}] Macro parsing and expansion, used during parsing.
    \item[\texttt{zmod.c}] Module resolution: handles \zn{use} imports and AST debug and logs printing.
    \item[\texttt{zmem.h}] Supports modular allocators.
    \item[\texttt{zmem.c}] Arena allocator for fast and linear allocations.
    \item[\texttt{zvec.h}] Header-only library for dynamic arrays.
    \item[\texttt{zhset}] Header-only library for hashset.
    \item[\texttt{zinc.h}] Header file that contains all definitions for the compilation pipeline.
\end{description}

\section{First Program}

\begin{lstlisting}[style=zinc]
libc :: foreign {
    printf :: u0 (*char, ...) // imports C functions
}

main :: i32 () {
    libc::printf("Hello world\n")
    return 0
}
\end{lstlisting}
